\section{The sharp support--margin law at fixed count budget}
\label{sec:frontier}

Fix \(a>0\), \(C\ge2\), and \(m\ge0\). For
\(0<\beta\le\min(a,1)\), let \(\mathcal Z(C,\beta;a,m)\)
be the class of divisors satisfying \eqref{eq:divisor}, containing
\(\beta\) and \(-\beta\), and obeying
\(N_0(Y)\le C(1+Y)^m\) for every \(Y\ge0\). Define
\[
 \cR_\eta(C,\beta;a,m)=\sup_{Z\in\mathcal Z(C,\beta;a,m)}r_Z(\eta).
\]
Other off-axis pairs are unrestricted except for the strip, reflection,
and count. The divisor, not only the test, varies in this supremum.

Put \(M=\lfloor C\rfloor\), \(p=2M-1\). The uniform conclusion is
\begin{equation}\label{frontier:statement}
 \cR_\eta(C,\beta;a,m)
 \asymp_{a,C,m}
 \beta^{-1}\log\bigl(1+(\beta\eta)^{1/p}\bigr),
 \qquad \eta\ge1,\quad0<\beta\le\min(a,1).
\end{equation}
Both comparison constants are positive and finite. They are uniform
in \(\beta,\eta,Z\), but no uniformity as \(C\) grows is
asserted. In particular,
\[
 \cR_1(C,\beta;a,m)\asymp_{a,C,m}\beta^{-(2M-2)/(2M-1)}.
\]
We will prove the upper bound by a smooth interpolation construction
and the lower bound against every test by a finite screen. A convenient
equivalent comparison function is \(\beta^{-1}\Phi_p(\beta\eta)\), where
\begin{equation}\label{frontier:comparison}
 \Phi_p(t)=
 \begin{cases}t^{1/p},&0<t\le1,\\ \log(1+t),&t>1.
 \end{cases}
\end{equation}
For fixed \(p\), this is comparable with
\(\log(1+t^{1/p})\) on the whole positive half-line.

\subsection{Three ordinate bands}
Choose \(0<\delta\le1\) so that
\(C(1+\delta)^m<M+1\). Such a choice exists by continuity;
if \(m=0\), take \(\delta=1\). Integer multiplicities imply
\(N_0(\delta)\le M\). List the distinct low locations as
\(z_1,\ldots,z_k\), so \(2\le k\le M\).

Choose any fixed \(R\ge16(a+2)\), and put
\[
 d_0=4/R,\qquad r_* =\lceil C(1+R)^m\rceil,\qquad V=a+R.
\]
Let \(v_1,\ldots,v_r\) be the distinct intermediate locations,
with \(\delta<|\Im v_j|\le R\); then \(r\le r_*\).
The intermediate annihilator and two useful bounds are
\begin{equation}\label{frontier:annihilator}
 A(w)=\prod_{j=1}^r(1-w/v_j),\qquad
 A_0=(1+\delta^{-1})^{r_*},\qquad
 A_1=(2V/\delta)^{r_*}.
\end{equation}
The coefficient sum of \(A\) is at most \(A_0\). If
\(|\Im w|\le\delta/2\), then
\(|1-w/v_j|>\delta/(2V)\), so \(|A(w)|^{-1}\le A_1\).
The empty product is one. Every actual intermediate location is
annihilated, irrespective of its multiplicity.

Fix an even nonnegative \(h\in C_c^\infty(-1,1)\) with
\(\int h=1\), and write \(G=H_h\). For \(|\Im s|\le1/4\),
\begin{equation}\label{frontier:bump-reciprocal}
 \Re G(s)\ge\cos(1/4)G(\Re s)\ge\cos(1/4)>1/2.
\end{equation}
The second inequality follows from evenness and \(\cosh x\ge1\).
Thus the reciprocal of \(G\) is holomorphic and bounded by two
in the narrower strip in which we will interpolate.

\subsection{Grouping close nodes without a separation assumption}
For now take free parameters
\begin{equation}\label{frontier:free}
 \ell\ge1,\qquad \epsilon=\beta\ell\le1,\qquad
 \tau\ge0,\qquad T=\tau\ell.
\end{equation}
Scale the low nodes to \(s_j=\ell z_j\), including
\(\epsilon,-\epsilon\), and set
\[
 \rho_0=\frac{\delta\epsilon}{64M(1+\tau)}.
\]
Choose \(\rho\in[\rho_0,2\rho_0]\) avoiding the finitely many
radii for which two boundary circles are tangent or three distinct
boundary circles meet. Form the union of the disks
\(|s-s_j|<\rho\), and let \(\Omega_+,\Omega_-\) be the
components containing \(\epsilon,-\epsilon\).

A path in the disk-intersection graph uses at most \(k-1\) edges
without repetition. Consequently every point of either selected
component's closure is within \(2k\rho\) of its target, and
\begin{equation}\label{frontier:thin}
 2k\rho\le\frac{\delta\epsilon}{16(1+\tau)},\qquad
 |\Im s|\le\frac{\delta\epsilon}{16(1+\tau)},\qquad |s|<2
 \quad(s\in\overline\Omega_+\cup\overline\Omega_-).
\end{equation}
The components are distinct, lie in opposite open half-planes, and
are interchanged by \(s\mapsto-\overline s\). Orient their
outer boundaries positively and any hole boundaries negatively.
They are finite piecewise smooth contours with total length at most
\(2\pi k\rho\). Every boundary point has distance at least
\(\rho\) from every low node. A node in a hole is therefore
excluded with winding number zero, as it should be.

Normalize the nodal polynomial by
\[
 p_0(s)=\prod_{j=1}^k\frac{s-s_j}{\sqrt{1+|s_j|^2}}.
\]
At a selected boundary point \(\zeta\), the normalized distance
to a node of modulus at most four is at least \(\rho/\sqrt{17}\).
For a larger node, \(|\zeta-s_j|>|s_j|/2\), giving a lower
bound \(1/\sqrt5\). Thus
\begin{equation}\label{frontier:nodal}
 |p_0(\zeta)|\ge(\rho/5)^k,\qquad
 \left\|\frac{p_0(\zeta)-p_0(s)}{\zeta-s}\right\|_{\coef,1}
 \le k3^{k-1}.
\end{equation}
Here \(\|\cdot\|_{\coef,1}\) denotes the sum of absolute
coefficients. To see the second estimate, telescope the difference
of the two products. A polynomial factor has coefficient sum at most
\(\sqrt2\), an evaluated factor has modulus at most \(\sqrt5\),
and the removed linear factor contributes at most one. Bounding the
remaining \(k-1\) factors by three proves the assertion. This is
the point at which normalizing far-away scaled nodes is useful.

Let \(b(w)=\sinh(d_0w)/(d_0w)\), with its value one at zero.
For real \(x,y\), elementary inequalities give
\[
 \left|\frac{\sinh(x+iy)}{x+iy}\right|^2
 \ge\left(\frac{\sin y}{y}\right)^2,\qquad
 \left|\frac{\sinh(x+iy)}{x+iy}\right|^{-K}
 \le e^{2Ky^2}\quad(|y|\le1/4).
\]
Indeed \(\sinh^2x\ge x^2\), \(\sin^2y\le y^2\), and
\(\sin y/y\ge1-y^2/6\); taking logarithms gives the second
bound. By \eqref{frontier:thin}, on the selected closures
\[
 |b(s/\ell)|^{-K}\le e^{\chi K},\qquad
 \chi=\frac{\delta^2}{8R^2}\frac{\beta^2}{(1+\tau)^2}.
\]
For any positive integer \(K\), the data
\[
 M_K(s)=\{A(s/\ell)b(s/\ell)^KG(s)\}^{-1}
\]
are holomorphic near both selected closures and bounded in modulus
by \(2A_1e^{\chi K}\). Nonvanishing holds on containing small
target disks, so possible holes cause no analytic difficulty.

The polynomial
\begin{equation}\label{frontier:interpolant}
 P(s)=\frac1{2\pi i}
 \left(\int_{\partial\Omega_+}-\int_{\partial\Omega_-}\right)
 M_K(\zeta)
 \frac{p_0(\zeta)-p_0(s)}{p_0(\zeta)(\zeta-s)}\,d\zeta
\end{equation}
has degree at most \(k-1\). Cauchy's formula gives the values
\(+M_K(s_j)\), \(-M_K(s_j)\), or zero, according as the node
is in \(\Omega_+\), \(\Omega_-\), or neither. The length
bound and \eqref{frontier:nodal} give
\begin{equation}\label{frontier:coefficients}
 \|P\|_{\coef,1}\le
 D_0\epsilon^{1-k}(1+\tau)^{k-1}e^{\chi K},\qquad
 D_0=2A_1M^2 3^{M-1}5^M(64M/\delta)^{M-1}.
\end{equation}
The boundary length supplies one power of \(\rho\). Losing that
power would give the wrong exponent in \eqref{frontier:statement}.
No inverse minimum separation has appeared.

\subsection{Selecting the tail budget before interpolation}
We now choose \(K\), rather than assume that a later choice can
absorb the size of the interpolant. All of the following are fixed
once \(a,C,m,R\) are fixed:
\[
 \begin{split}
 D_{\rm tail}&=D_0(2R)^{M-1}(3R/\delta)^{r_*},\\
 S_R&=2^{m+1}C(1+R)^m,\qquad B_R=\tfrac12\log(2S_R),\\
 d_*&=M+r_*+m,\qquad J_*=1+d_*+\log D_{\rm tail}+B_R.
 \end{split}
\]
Put \(u_\beta=\log(1/\beta)\), and select
\begin{equation}\label{frontier:K}
 \begin{split}
 E&=a\ell(1+\tau)+(M-1)\{u_\beta+\log(1+\tau)\}
                      +\log D_{\rm tail},\\
 K&=\left\lceil\max\{d_*,E+B_R\}\right\rceil
 \le a\ell(1+\tau)+(M-1)\{u_\beta+\log(1+\tau)\}+J_*.
 \end{split}
\end{equation}
Use this integer in \eqref{frontier:interpolant}. The complete
transform will be
\begin{equation}\label{frontier:transform}
 \mathcal F(w)=\sqrt\ell\,P(\ell w)G(\ell w)A(w)b(w)^K\cosh(Tw).
\end{equation}
Before the final \(\cosh\) factor, its low-node values are
\(+\sqrt\ell,-\sqrt\ell,0\), as prescribed, and all its
intermediate values vanish.

For \(w=\sigma+it\), \(|\sigma|\le a\), \(v=|t|>R\),
we have
\[
 \begin{aligned}
 1+\ell|w|&\le2\ell v,&
 |A(w)|&\le(3R/\delta)^{r_*}(v/R)^r,\\
 |G(\ell w)|&\le e^{a\ell},&
 |\cosh(Tw)|&\le e^{a\tau\ell},\\
 |b(w)|&\le e^{4a/R}\frac R{4v}.&&
 \end{aligned}
\]
In the polynomial bound the factors of \(\ell\) cancel exactly:
\(\epsilon^{1-k}(2\ell v)^{k-1}=\beta^{1-k}(2v)^{k-1}\).
Moreover \(\chi\le1/16\), \(4a/R<1/4\), and hence
\(\log4-4a/R-\chi>1\). Substituting
\eqref{frontier:coefficients} into \eqref{frontier:transform}
therefore gives
\begin{equation}\label{frontier:tailpoint}
 |\mathcal F(\sigma+it)|\le
 \sqrt\ell\,e^{E-K}(v/R)^{-q_*},\qquad
 q_*=K-(k-1)-r\ge m+1.
\end{equation}
Dyadic summation gives
\(\sum_{|\Im z|>R}\nu(z)(|\Im z|/R)^{-2q_*}\le S_R\).
Apply \eqref{frontier:tailpoint} at both reflected locations to obtain
\begin{equation}\label{frontier:tail}
 \sum_{|\Im z|>R}\nu(z)
  |\mathcal F(z)\overline{\mathcal F(z^\#)}|
 \le\ell e^{2(E-K)}S_R\le\ell/2.
\end{equation}
The estimate concerns the full reflected tail, not a diagonal
substitute for it.

\subsection{Signs, norm, and physical support}
Reflection exchanges the selected components. A selected pair at
\(z\) contributes
\(-2\nu(z)\ell\Re\cosh^2(Tz)\). By
\eqref{frontier:thin}, \(|T\Im z|\le1/16\). The identity
\[
 \Re\cosh^2(x+iy)=\tfrac12\{1+\cosh(2x)\cos(2y)\}
\]
shows that every selected pair contributes negatively. Keeping the
exact contribution of the targets until after tail absorption gives
\begin{equation}\label{frontier:negative}
 Q_Z(f,f)\le-2\ell\cosh^2(\tau\epsilon)+\ell/2
                 \le-\tfrac12\ell e^{2\tau\epsilon}.
\end{equation}
All other low-node contributions and all intermediate contributions
are zero.

Here the inverse transform is an actual smooth test. Let \(\mu\)
be uniform probability measure on \([-d_0,d_0]\), and put
\[
 h_\ell(u)=\ell^{-1/2}h(u/\ell),\qquad
 f_0=\mu^{*K}*\{A(-\partial_u)P(-\ell\partial_u)h_\ell\},\qquad
 f(u)=\tfrac12\{f_0(u-T)+f_0(u+T)\}.
\]
Integration by parts and Fubini give \(H_f=\mathcal F\).
Differentiation preserves the initial bump's strict support gap, and
convolution by a compact measure preserves smoothness. Thus
\begin{equation}\label{frontier:support-exact}
 f\in C_c^\infty(-L_{\rm out},L_{\rm out}),\qquad
 L_{\rm out}=\ell(1+\tau)+4K/R.
\end{equation}

The order of operations yields an efficient norm estimate. A term
of degree \(j\) from \(P\) and degree \(b\) from \(A\)
acts on the bump as a scalar multiple of
\(\ell^{-b}\ell^{-1/2}h^{(j+b)}(u/\ell)\). Its norm is at
most \(\|h^{(j+b)}\|_2\), since \(\ell\ge1\). Probability
convolution and the average of two translations are \(L^2\)
contractions. Hence, with
\(H_* =\max_{0\le j\le M-1+r_*}\|h^{(j)}\|_2\),
\begin{equation}\label{frontier:norm}
 \norm f_2\le A_0H_*\|P\|_{\coef,1}.
\end{equation}
In particular, the convolution incurs no exponential norm cost in
\(K\), and the intermediate differential operator incurs no
positive power of \(\ell\).

Write \(\mu_0=\delta^2/(8R^2)\). Multiplying
\eqref{frontier:K} by \(\chi=\mu_0\beta^2/(1+\tau)^2\),
and using \(\beta\ell\le1\), gives
\[
 \chi K\le\mu_0\{a+(M-1)/e+J_*\}=:H_{\rm inv}.
\]
Here we used \(\beta^2\log(1/\beta)\le1/(2e)\) and
\(\log(1+\tau)/(1+\tau)^2\le1/(2e)\).
Let
\[
 B=\max\{1,(A_0H_*D_0)^2e^{2H_{\rm inv}}\}.
\]
Equations \eqref{frontier:coefficients} and \eqref{frontier:norm}
then give
\[
 \|f\|_2^2\le B\left(\frac{1+\tau}{\epsilon}\right)^{2M-2}.
\]
The negative value proves \(f\ne0\). We have proved the following
construction for every set of parameters in \eqref{frontier:free}:
\begin{equation}\label{frontier:master}
 \begin{split}
 -\frac{Q_Z(f,f)}{\norm f_2^2}
 &\ge c\ell\left(\frac{\beta\ell}{1+\tau}\right)^{p-1}
                      e^{2\tau\beta\ell},\\
 L_{\rm out}&\le D\{\ell(1+\tau)+\log(1/\beta)+1\},
 \end{split}
\end{equation}
where \(c=1/(2B)>0\) and a finite \(D\) depend only on
\(a,C,m,R,h\). For the second line, use
\(\log(1+\tau)\le\tau\), \(\ell\ge1\) in
\eqref{frontier:K}. For example
\(D=1+4a/R+4(M-1)/R+4J_*/R\) suffices. We have retained
\eqref{frontier:K} and \eqref{frontier:support-exact} separately
because their leading dependence on \(T\) will yield an exact
growth exponent later.

\subsection{The power branch and the logarithmic branch}
For \eqref{frontier:statement} fix \(R=16(a+2)\). Let
\(t=\beta\eta\), \(q=(p-1)/p\), choose \(A\ge1\) with
\(cA^p\ge2\), and put \(t_0=A^{-p}\).
If \(0<t\le t_0\), set
\[
 \ell=A\beta^{-1}t^{1/p}=A\eta^{1/p}\beta^{-q},\qquad\tau=0.
\]
Then \(\ell\ge1\), \(\beta\ell\le1\), and the margin in
\eqref{frontier:master} is at least \(2\eta\).
The elementary maximum of \(x^q\log(1/x)\) gives
\(\log(1/\beta)\le\beta^{-q}/(qe)\le\ell/(qe)\).
Since \(q\ge2/3\), the support is at most
\(3DA\beta^{-1}t^{1/p}\).

For \(t>t_0\), take \(\ell=1/\beta\), so \(\epsilon=1\).
Choose
\[
 H=\max\{1,4(p-1)^2,2\log(2/c)\},\qquad
 K_1=\max\{1,H/\log2\},\qquad \tau=K_1\log(2+t).
\]
For \(\tau\ge H\), the inequality
\(\log(1+\tau)\le\sqrt\tau\) implies
\((p-1)\log(1+\tau)\le\tau/2\). One direct verification
of the logarithmic inequality differentiates
\(\sqrt x-\log(1+x)\), whose derivative is
\((\sqrt x-1)^2/(2\sqrt x(1+x))\ge0\).
It follows that \(ce^\tau/(1+\tau)^{p-1}\ge2\).
The margin is therefore at least
\(2\beta^{-1}e^\tau>2\eta\), while the support is bounded by
a fixed constant times \(\beta^{-1}\log(2+t)\).
For \(t\ge1\), this is at most twice the logarithmic scale in
\eqref{frontier:comparison}. For the remaining interval
\(t_0<t\le1\), use
\(\log(2+t)\le\log3\le A\log3\,t^{1/p}\).
This covers the entire transition interval and proves the upper
bound, always with a strict margin reserve.

\subsection{A matching obstruction for every test}
Let \(n=M-2\), so \(p=2n+3\), and define
\begin{equation}\label{frontier:lower-screen}
 T_n^2=(n+1)^3 4^n,\qquad R_n=24T_n^2,\qquad
 Z_{n,\beta}=[\beta]+[-\beta]+\sum_{j=1}^n[i\beta R_nj].
\end{equation}
The notation \([z]\) denotes one atom. This divisor has total mass
\(M\le C\), and so satisfies the whole polynomial counting bound.

For every polynomial \(P\) of degree at most \(n\), with
coefficient vector \(b\), we first claim
\begin{equation}\label{frontier:moment-positive}
 2\Re(P(1)\overline{P(-1)})+\sum_{j=1}^n|P(iR_nj)|^2
 \ge(2T_n^2)^{-1}\|b\|_2^2.
\end{equation}
Lagrange interpolation at \(0,i,\ldots,in\) gives basis
polynomials of coefficient sum at most \((n+1)2^n\): the
denominator at \(ij\) has modulus \(j!(n-j)!\), while the
numerator's coefficient sum is at most \((n+1)!\).
Cauchy--Schwarz consequently gives
\[
 2|p(0)|^2+\sum_{j=1}^n|p(ij)|^2\ge T_n^{-2}\|\coef(p)\|_2^2.
\]
For \(0\le h\le1/2\), a geometric-series estimate gives
\(|p(\pm h)-p(0)|\le2h\|\coef(p)\|_2\).
Replacing \(2|p(0)|^2\) by
\(2\Re(p(h)\overline{p(-h)})\) changes the form by at most
\(12h\|\coef(p)\|_2^2\). Take \(h=1/R_n\) and apply this
to \(p(w)=P(R_nw)\); its coefficient norm is at least that
of \(P\). This proves \eqref{frontier:moment-positive}, including
\(n=0\).

For \(g\in L^2(-1,1)\), Taylor expand
\[
 H_g(\epsilon s)=P_g(s)+r_g(s),\qquad
 P_g(s)=\sum_{k=0}^n\frac{\epsilon^ks^k}{k!}\int_{-1}^1x^kg(x)\,dx.
\]
Set \(K_0=\max(1,nR_n)\). At each of the fixed nodes
\(1,-1,iR_n,\ldots,inR_n\), the exponential remainder gives
\[
 |r_g(s)|\le E_n\epsilon^{n+1}\|g\|_2,\qquad
 E_n=\frac{\sqrt2 eK_0^{n+1}}{(n+1)!},\qquad K_0\epsilon\le1.
\]
Let \(V_n\) evaluate coefficient vectors at these nodes and let
\(J_0\) swap the first two coordinates and fix the others.
Its norm is one. With \(b=\coef(P_g)\), the polynomial part
is at least \(b_n\|b\|^2\), where \(b_n=(2T_n^2)^{-1}\).
The cross term is bounded by \(2\|V_n\|\|b\|\|r\|\).
Completing the square gives
\[
 (V_nb+r)^*J_0(V_nb+r)
 \ge-(1+\|V_n\|^2/b_n)\|r\|^2
 \ge-D_n\epsilon^{2n+2}\|g\|_2^2,
\]
with \(D_n=(1+\|V_n\|^2/b_n)(n+2)E_n^2\).
Scaling from \((-1,1)\) to \((-L,L)\) yields the all-test estimate
\begin{equation}\label{frontier:all-test-lower}
 Q_{Z_{n,\beta}}(f,f)\ge-D_nL(\beta L)^{p-1}\|f\|_2^2
 \qquad(\beta L\le K_0^{-1}).
\end{equation}
This bounds the original form, not merely a polynomial compression.

For \(t=\beta\eta\le1\), choose \(k_1>0\) so that
\(k_1\le K_0^{-1}\), \(D_nk_1^p\le1/2\), and take
\(L=k_1\beta^{-1}t^{1/p}\). The envelope is at most
\(\eta/2\), excluding the requested margin at this radius and
all smaller radii. For \(t\ge1\), axis positivity and
\eqref{eq:single-pair} instead imply
\[
 \cM_{Z_{n,\beta}}(L)\le e^{2\beta L}/(2\beta).
\]
At \(L=(4\beta)^{-1}\log(1+t)\), its ratio to \(\eta\)
is at most \(\sqrt{1+t}/(2t)\le\sqrt2/2<1\).
Taking the smaller lower constant proves
\eqref{frontier:statement} in both regimes with the same finite screen.

\subsection{The exact role of centered conjugation symmetry}
If in addition \(Z=\overline Z\) in the coordinates used above,
let \(\cR_\eta^{\rm sym}\) denote the supremum over this smaller
class. Then
\begin{equation}\label{frontier:symmetric}
 \cR_\eta^{\rm sym}(C,\beta;a,m)
 \asymp_{a,C,m}
 \beta^{-1}\log\bigl(1+(\beta\eta)^{1/p_e}\bigr),\qquad
 p_e=4\lfloor C/2\rfloor-1.
\end{equation}
Here too \(\eta\ge1\) and \(0<\beta\le\min(a,1)\).

To obtain the upper bound, omit zero from the low interpolation list.
Every remaining location has its distinct negative with equal
multiplicity, so the number of distinct nonzero low nodes is even
and at most \(M_e=2\lfloor C/2\rfloor\). Use \(M_e\) in all
local degree and contour budgets, retaining the original count for
the intermediate and tail bounds. The intermediate polynomial
\(A\) is even and real, as are \(G\) and \(b\).
The selected positive component is conjugation-invariant and the
negative component is its negative. The prescribed interpolation
data are odd and conjugation-compatible. Uniqueness of polynomial
interpolation therefore shows that \(P\) is odd and has real
coefficients. The resulting physical test is real and odd. Its
transform vanishes at zero, accounting for every omitted zero
multiplicity. All estimates already proved hold with exponent
\(2M_e-1=p_e\).

For the lower bound write \(M_e=2r+2\), and use
\[
 [\beta]+[-\beta]+\sum_{j=1}^r
        \{[i\beta R_rj]+[-i\beta R_rj]\},\qquad
 T_r^2=(2r+1)(r+1)^{4r},\qquad R_r=24T_r^2.
\]
Lagrange interpolation on \(-ir,\ldots,0,\ldots,ir\)
has coefficient bound \(T_r\) by the same argument: each
numerator coefficient sum is at most \((r+1)^{2r}\), and
each denominator modulus is at least one. The perturbation from
the repeated zero evaluations to \(\pm1/R_r\) costs at most
\(12/R_r\) times the squared coefficient norm. It follows that
the signed evaluation form is bounded below by
\((2T_r^2)^{-1}\|\coef(P)\|^2\) on all complex polynomials
of degree at most \(2r\). Taylor truncation at that degree and
the identical square completion give
\[
 Q_Z(f,f)\ge-D_rL(\beta L)^{4r+2}\|f\|_2^2
 \quad\bigl(\beta L\le\max(1,rR_r)^{-1}\bigr)
\]
for a finite constant \(D_r\). Since \(4r+3=p_e\), this is
the required small-margin obstruction; \eqref{eq:single-pair}
gives the logarithmic obstruction. Both hold for all complex tests,
not only odd ones.

The qualification ``centered'' is essential. If an original divisor
is globally conjugation-symmetric and is moved by \(-i\gamma\)
to put a target at ordinate zero, its original conjugation becomes
\(w\mapsto\overline w-2i\gamma\), not
\(w\mapsto\overline w\). Thus ordinary global conjugation does
not justify \eqref{frontier:symmetric} at a nonzero target height.

\subsection{Fixed continuous counting profiles}
The proof uses only an integer low-node budget, a finite intermediate
budget, and a polynomial tail. Consequently, for each fixed continuous
nondecreasing \(g:[1,\infty)\to[1,\infty)\), with \(g(1)=1\)
and \(g(2t)\le\Delta g(t)\), both laws hold under
\(N_0(Y)\le C(1+Y)g(1+Y)\), with constants depending on
\(a,C,g\). Indeed, continuity supplies
\(C(1+\delta)g(1+\delta)<\lfloor C\rfloor+1\), and
doubling gives \(g(t)\le\Delta t^{\log_2\Delta}\).
The finite lower examples still satisfy the entire count because
their total mass is at most \(C\). A common doubling constant
does not supply a common modulus of continuity at one, so this
observation asserts no uniformity over all such profiles.
